\documentclass[12pt,a4paper,oneside]{book} % twoside for draf

%\usepackage{babel}
\usepackage[utf8]{vietnam}
%\usepackage{times}
\usepackage{graphicx}
\usepackage{subcaption}

\usepackage{mathptmx}	% same Time New Roma
%\renewcommand{\rmdefault}{phv} % Arial
%\renewcommand{\sfdefault}{phv} % Arial
\usepackage{url}

%link to part by click on contents region
\usepackage{hyperref}
\hypersetup{
    colorlinks,
    citecolor=black,
    filecolor=black,
    linkcolor=black,
    urlcolor=black
}

\usepackage{fancyhdr}
\usepackage{algorithm2e}
\usepackage{amsmath}
\usepackage{array,tabularx}
\usepackage{amsfonts}
\usepackage{amssymb}
\usepackage{cases}
\usepackage{tabularx}
\usepackage{enumitem}

\usepackage{xurl}
\urlstyle{tt}
\DeclareUrlCommand\codepath{\urlstyle{tt}\Urlmuskip=0mu plus 1mu\relax}
\interfootnotelinepenalty=10000
\usepackage{adjustbox}
\usepackage{multirow}
\usepackage{enumitem}
\usepackage{bkthesis}
\usepackage{tikz}
\usetikzlibrary{shapes,arrows}

\usepackage{acronym} 
\usepackage{float}
\usepackage{algpseudocode}
\usepackage{tikz}
\usepackage{pgfplots}
\usepackage{pdfpages}
\usepackage{svg}
\svgsetup{inkscapelatex=true}
\pgfplotsset{compat=1.17}
\newenvironment{conditions*}
  {\par\vspace{\abovedisplayskip}\noindent
   \tabularx{\columnwidth}{>{$}l<{$} @{${}:{}$} >{\raggedright\arraybackslash}X}}
  {\endtabularx\par\vspace{\belowdisplayskip}}

\tikzstyle{blank} = [text badly centered, node distance = 2cm, inner sep=0pt]
\tikzstyle{block} = [rectangle, draw, 
    text width=2em, text centered, minimum height=2em]
\tikzstyle{line} = [draw, -latex']

\newcolumntype{b}{>{\hsize=1.0\hsize}X}

\crname{BÁO CÁO\\THỰC TẬP 1}
\ctname{KHẢO SÁT VÀ ĐÁNH GIÁ CÁC PHƯƠNG PHÁP PHÁT HIỆN GIỌNG NÓI GIẢ MẠO}
\cmajorName{Ngành: Khoa học Máy tính}
\cstuname{
	Học viên thực hiện: \\
	\begin{tabular}{ l c }
		Bùi Hoàng Minh & MSHV: 2470504
  % \\
		% DEF & MSSV: 15xxxxx
	\end{tabular}
}

% \csCouncil{Hội đồng 5 Khoa học Máy Tính}
\csSupervise{TS. Nguyễn Đức Dũng}
% \csPhanbien{ThS. Võ Thanh Hùng}
\csReviewer{}
\cttime{05/2026}

\thesislayout

\begin{document}
%-	Bìa cứng - màu xanh dương, chữ mạ vàng (xem mẫu đính kèm)
%-	Trang tên (tờ lót): chất liệu giấy, nội dung giống như bìa LV
%-	Ở gáy LV: in nhan đề LV (có thể in tóm tắt nếu nhan đề quá dài), size 15 – 17
%-	Phiếu Nhiệm vụ LV, chấm điểm Hướng dẫn & Phản biện (đã ký): nhận từ GVHD & GVPB sau khi bảo vệ (theo lịch hẹn).
%-	Lời cam đoan
%-	Lời cảm ơn/ Lời ngỏ
%-	Tóm tắt LV
%-	Mục lục
%-	Danh mục, bảng biểu, hình ảnh, ... (nếu có)
%-	Nội dung LV
%-	Danh mục TL tham khảo
%-	Phụ lục (nếu có)

\coverpage

\frontmatter

% add content here
%-	Lời cam đoan
\begin{declaration}
Tôi xin cam đoan đây là công trình nghiên cứu của riêng tôi dưới sự hướng dẫn của TS. Nguyễn Đức Dũng. Nội dung nghiên cứu và các kết quả đều là trung thực và chưa từng được công bố trước đây. Các số liệu được sử dụng cho quá trình phân tích, nhận xét được chính chúng tôi thu thập từ nhiều nguồn khác nhau và sẽ được ghi rõ trong phần tài liệu tham khảo.

Ngoài ra, chúng tôi cũng có sử dụng một số nhận xét, đánh giá và số liệu của các tác giả khác, cơ quan tổ chức khác. Tất cả  đều có trích dẫn và chú thích nguồn gốc.

Nếu phát hiện có bất kì sự gian lận nào, chúng tôi xin hoàn toàn chịu trách nhiệm về nội dung thực tập của mình. Trường đại học Bách Khoa thành phố Hồ Chí Minh không liên quan đến những vi phạm tác quyền, bản quyền do chúng tôi gây ra trong quá trình thực hiện.

\end{declaration}
~

\begin{acknowledgments}
Trong quá trình nghiên cứu đề tài Thực tập 1, nhóm đã nhận được nhiều sự giúp đỡ và đóng góp quý báu. Đầu tiên, nhóm xin gửi lời cảm ơn sâu sắc đến TS. Nguyễn Đức Dũng, người đã tận tâm chia sẻ kinh nghiệm và kiến thức, đồng thời hỗ trợ nhóm hết mình trong thời gian hướng dẫn nhóm thực hiện đồ án môn học Thực tập 1. Tiếp theo, nhóm xin chân thành cảm ơn các thầy, cô giáo trong trường Đại học Bách Khoa Hồ Chí Minh nói chung và khoa Khoa học và Kỹ thuật Máy tính nói riêng đã luôn giúp đỡ, truyền đạt kiến thức và tạo điều kiện mọi mặt trong thời gian nhóm học tập tại trường. Cuối cùng, nhóm xin cảm ơn gia đình và bạn bè, những người đã quan tâm đến tiến độ của đề tài và tạo động lực để nhóm hoàn thành Thực tập 1.
	
\end{acknowledgments}
~

%
\begin{arc}
\begin{acronym}[ASVspoof]
    \acro{SDD}{Speech Deepfake Detection}
    \acro{TTS}{Text-to-Speech}
    \acro{VC}{Voice Conversion}
    \acro{CM}{Countermeasure}
    \acro{ASV}{Automatic Speaker Verification}
    \acro{EER}{Equal Error Rate}
    \acro{FAR}{False Acceptance Rate}
    \acro{FRR}{False Rejection Rate}
    \acro{LA}{Logical Access}
    \acro{PA}{Physical Access}
    \acro{SSL}{Self-Supervised Learning}
    \acro{LFCC}{Linear Frequency Cepstral Coefficients}
    \acro{CQCC}{Constant-Q Cepstral Coefficients}
    \acro{LCNN}{Light Convolutional Neural Network}
    \acro{ROC}{Receiver Operating Characteristic}
    \acro{DET}{Detection Error Tradeoff}
\end{acronym}
\end{arc}


\begin{abstract}

Báo cáo này khảo sát và đánh giá các phương pháp phát hiện giọng nói giả mạo trong bối cảnh các kỹ thuật tổng hợp giọng nói và chuyển đổi giọng nói ngày càng dễ tiếp cận. Trọng tâm của báo cáo là bài toán Speech Deepfake Detection (SDD), trong đó hệ thống cần phân biệt một đoạn audio đầu vào là giọng thật (bonafide) hay giọng giả mạo (spoof). Báo cáo trình bày cơ sở lý thuyết về các dạng tấn công, các tập dữ liệu benchmark, các nhóm kiến trúc mô hình và phương pháp đánh giá.

Bên cạnh phần khảo sát, báo cáo tái lập benchmark chính với năm mô hình tiêu biểu gồm LFCC+LCNN, AASIST, AASIST-L, XLS-R+AASIST và XLS-R+Nes2Net trên bốn tập dữ liệu: ASVspoof 2019 LA, ASVspoof 2021 DF, ASVspoof 5 Track 1 và In-the-Wild. AASIST3 được khảo sát ở mức checkpoint công khai nhưng không đưa vào phân tích chính vì model card Hugging Face cho biết checkpoint hiện tại đã deprecated và không phải weights dùng trong kết quả paper. Kết quả thực nghiệm cho thấy hiệu năng của các mô hình thay đổi đáng kể khi chuyển từ benchmark sạch sang dữ liệu có codec, tấn công hiện đại hoặc điều kiện in-the-wild. Nhóm mô hình sử dụng front-end self supervised learning (SSL), đặc biệt XLS-R+Nes2Net, đạt EER thấp hơn và generalization gap nhỏ hơn trên hầu hết các dataset được khảo sát. Phân tích lỗi trên ASVspoof 5 chỉ ra hai failure mode chính: codec dominance (EER nhóm XLS-R tăng từ 2.71-3.78\% trên nocodec lên 35-42\% trên C04/C07) và confident false positive theo attack ID (A28 đứng đầu FP rate cho cả hai mô hình XLS-R, A17 đứng đầu cho AASIST), tạo động lực cho hai hướng ứng viên cụ thể là codec-aware training và attack-aware hard mining. Tuy nhiên, phạm vi hiện tại vẫn là một pilot benchmark. Vì vậy, báo cáo đề xuất hướng phát triển tiếp theo theo thứ tự: mở rộng bản đồ dữ liệu và kiến trúc cho SDD, chuẩn hoá protocol đánh giá, sau đó chọn baseline và kiểm chứng các hướng cải thiện ở mức ý niệm.
\end{abstract}
 

% sử dụng các phần bên dưới nếu cần:
\tableofcontents
%\listofsymbols
\listoftables
\listoffigures
%\listofalgorithms


\mainmatter

\fancypagestyle{plain}{%
\fancyhf{}
\fancyfoot[LE, RO]{\thepage}
\renewcommand{\headrulewidth}{0.0pt}
}

\pagestyle{fancy}
\renewcommand{\chaptermark}[1]{\markboth{\MakeUppercase{#1}}{}}
\fancyhf{}
\fancyhead[LE,RO]{\leftmark}
\fancyhead[RE,LO]{}
\fancyfoot[LE,RO]{\thepage}
\renewcommand{\footrulewidth}{0.4pt}

%\fancyhead{}  % Clears all page headers and footers
%\rhead{\thepage}  % Sets the right side header to show the page number
%\lhead{}  % Clears the left side page header
%\fancyfoot[positions]{footer}
%\renewcommand{\footrulewidth}{0.4pt}

%\pagestyle{fancy}  % Finally, use the "fancy" page style to implement the FancyHdr headers

% \input{chapters/translation-table.tex}

\chapter{Giới thiệu} \label{chap:intro}

\section{Bối cảnh và động lực}

Trong khoảng một thập kỷ trở lại đây, các kỹ thuật tổng hợp giọng nói (Text-to-Speech, TTS) và chuyển đổi giọng nói (Voice Conversion, VC) đã có những bước tiến vượt bậc. Từ các mô hình tự hồi tiếp như WaveNet và Tacotron đến các kiến trúc end-to-end hiện đại như VITS, neural codec TTS, hay các mô hình dựa trên diffusion, chất lượng giọng nói tổng hợp ngày càng tiệm cận với giọng nói thật về cả độ tự nhiên lẫn độ giống người nói gốc. Nhiều hệ thống thương mại hiện nay có thể nhân bản giọng nói từ chỉ vài giây mẫu âm thanh, mở ra nhiều ứng dụng tích cực như trợ lý ảo, đọc sách nói, hay khôi phục giọng cho người mất khả năng nói.

Tuy nhiên, sự phát triển nhanh của các kỹ thuật này cũng kéo theo những rủi ro nghiêm trọng. Các vụ lừa đảo qua điện thoại bằng giọng nói giả mạo người thân, mạo danh lãnh đạo doanh nghiệp để lừa chuyển khoản, hay phát tán thông tin sai lệch bằng cách giả giọng nhân vật công chúng đã được ghi nhận tại nhiều quốc gia trong những năm gần đây. Rủi ro này càng đáng quan ngại khi các công cụ voice cloning ngày càng dễ tiếp cận và chi phí thấp.

Trước thực trạng đó, bài toán phát hiện giọng nói giả mạo (Speech Deepfake Detection, SDD) trở thành một hướng nghiên cứu quan trọng. Ở dạng cơ bản nhất, bài toán được phát biểu như một tác vụ phân loại nhị phân: cho một đoạn audio đầu vào, hệ thống cần xác định đoạn đó là giọng thật (\textit{bonafide}) hay đã được tổng hợp/chỉnh sửa (\textit{spoof}). Đầu ra của hệ thống thường là một điểm số liên tục phản ánh xác suất \textit{bonafide}, từ đó có thể đưa ra quyết định phân loại bằng cách so sánh với một ngưỡng. Để đánh giá hệ thống, cộng đồng nghiên cứu thường sử dụng chỉ số Equal Error Rate (EER) -- tỉ lệ lỗi tại điểm cân bằng giữa False Acceptance Rate và False Rejection Rate, càng thấp càng tốt.

\section{Các thách thức chính}

Mặc dù đã có nhiều tiến bộ, các hệ thống SDD hiện nay vẫn đối mặt với một số thách thức cốt lõi.

\textbf{Khả năng tổng quát hoá liên miền.} Nhiều mô hình được báo cáo có hiệu năng cao trên benchmark mà chúng được huấn luyện, ví dụ ASVspoof 2019 LA, nhưng EER tăng đáng kể khi được đánh giá trên dữ liệu có phân bố khác như audio đã qua nén lossy trong ASVspoof 2021 DF hoặc dữ liệu được thu thập từ mạng xã hội trong In-the-Wild. Hiện tượng này cho thấy nhiều mô hình đang học các đặc trưng phụ thuộc vào điều kiện thu âm hoặc artifact cụ thể của tập huấn luyện, thay vì các dấu hiệu tổng quát của giọng nói giả mạo.

\textbf{Tốc độ tiến hoá của các kỹ thuật tấn công.} Các phương pháp TTS/VC mới liên tục xuất hiện, đặc biệt là các kiến trúc dựa trên neural codec, diffusion và adversarial training. Nhiều dataset benchmark được xây dựng cách đây vài năm có thể không bao gồm các loại tấn công mới nhất, dẫn đến rủi ro mô hình không nhận biết được các dạng tấn công mà nó chưa từng thấy. Việc đảm bảo các benchmark cập nhật kịp thời là một thách thức về mặt dữ liệu và quy trình đánh giá.

\textbf{Điều kiện thực tế.} Audio trong thực tế thường đi kèm các yếu tố như nén lossy (MP3, AAC, OPUS), nhiễu nền, microphone chất lượng kém hoặc qua nhiều khâu xử lý hậu kỳ. Các yếu tố này có thể che lấp hoặc làm biến dạng các dấu vết phổ mà nhiều mô hình dựa vào để phát hiện giọng nói giả mạo. Vì vậy, các mô hình huấn luyện trên audio sạch trong phòng thí nghiệm thường suy giảm hiệu năng đáng kể khi áp dụng cho audio thực tế.

\textbf{Mất cân bằng và đa dạng của dữ liệu.} Trong nhiều dataset, số lượng mẫu \textit{spoof} lớn hơn nhiều so với \textit{bonafide}; chẳng hạn ASVspoof 2019 LA có khoảng 7,3 nghìn mẫu \textit{bonafide} trên 64 nghìn mẫu \textit{spoof}. Đồng thời, độ đa dạng của \textit{bonafide} về người nói, ngôn ngữ và điều kiện thu âm thường hạn chế hơn. Điều này có thể ảnh hưởng đến khả năng học các đặc trưng tổng quát của giọng nói thật, đặc biệt khi áp dụng sang các tập dữ liệu cross-domain.

\section{Hướng tiếp cận và cấu trúc báo cáo}

Báo cáo này được thực hiện như bước khởi đầu trong một quy trình nghiên cứu dài hơi gồm Thực tập 2 và Luận văn tốt nghiệp. Phạm vi của Thực tập 1 tập trung vào ba hoạt động chính: khảo sát cơ sở lý thuyết về dữ liệu, kiến trúc và phương pháp đánh giá; tái lập một số mô hình tiêu biểu trên các tập dữ liệu phổ biến để có số liệu tham chiếu thực tế; và phân tích lỗi sơ bộ trên một dataset trọng tâm, từ đó đưa ra đề xuất hướng tiếp cận ở mức ý niệm. Các đề xuất chi tiết về thiết kế thí nghiệm và triển khai thực tế sẽ được phát triển ở các giai đoạn tiếp theo.

Cụ thể, báo cáo tái lập benchmark chính với năm mô hình đại diện cho các nhóm tiếp cận khác nhau (LFCC+LCNN, AASIST, AASIST-L, XLS-R+AASIST, XLS-R+Nes2Net) trên bốn dataset phổ biến: ASVspoof 2019 LA, ASVspoof 2021 DF, ASVspoof 5 và In-the-Wild. Toàn bộ thí nghiệm được thực hiện trên nền tảng Kaggle với GPU T4, chủ yếu sử dụng checkpoint công bố công khai; riêng LFCC+LCNN được huấn luyện lại như baseline do không có checkpoint phù hợp. AASIST3 được chạy thử từ checkpoint Hugging Face công khai nhưng bị loại khỏi benchmark chính vì checkpoint này không đại diện cho weights dùng trong paper. Phân tích lỗi sơ bộ được tập trung vào ASVspoof 5 -- tập dữ liệu mới nhất trong chuỗi ASVspoof Challenge với Track 1 eval gồm 16 attack và metadata codec phong phú -- nhằm đưa ra cái nhìn cụ thể về điểm mạnh và điểm yếu của từng nhóm mô hình trong các kịch bản tấn công gần đây.

\begin{figure}[H]
    \centering
    \begin{adjustbox}{width=\textwidth}
    \begin{tabular}{c c c c c c c}
        \fbox{\begin{minipage}{3.0cm}\centering Audio đầu vào\\\textit{bonafide}/\textit{spoof}\end{minipage}}
        & $\longrightarrow$ &
        \fbox{\begin{minipage}{3.0cm}\centering Front-end\\LFCC, waveform, SSL\end{minipage}}
        & $\longrightarrow$ &
        \fbox{\begin{minipage}{3.0cm}\centering Back-end\\LCNN, AASIST, Nes2Net\end{minipage}}
        & $\longrightarrow$ &
        \fbox{\begin{minipage}{3.0cm}\centering Score\\và quyết định\end{minipage}}
    \end{tabular}
    \end{adjustbox}
    \caption{Quy trình tổng quát của hệ thống phát hiện giọng nói giả mạo.}
    \label{fig:sdd-pipeline}
\end{figure}

Phần còn lại của báo cáo được tổ chức như sau. Chương \ref{chap:background} trình bày cơ sở lý thuyết, bao gồm phát biểu bài toán, các loại tấn công, các tập dữ liệu benchmark, các nhóm kiến trúc tiếp cận và phương pháp đánh giá. Chương \ref{chap:experiment} mô tả thiết lập thí nghiệm, trình bày kết quả EER tổng hợp trên năm mô hình trong benchmark chính và bốn dataset, đồng thời giải thích trường hợp checkpoint AASIST3 bị loại khỏi phân tích chính. Chương \ref{chap:conclusion} tổng kết các phát hiện chính, đưa ra đề xuất hướng tiếp cận ở mức ý niệm và các bước phát triển tiếp theo trong khuôn khổ Thực tập 2 và Luận văn.

\chapter{Cơ sở lý thuyết} \label{chap:background}

\section{Bài toán và các loại tấn công}

Speech Deepfake Detection (SDD) là bài toán xây dựng một hệ thống countermeasure (CM) để phân biệt audio thật (\textit{bonafide}) và audio giả mạo (\textit{spoof}). Ở mức utterance-level, với một đoạn tín hiệu âm thanh đầu vào $x$, mô hình sinh ra một điểm số:
\begin{equation}
    s = f(x) \in \mathbb{R},
\end{equation}
thể hiện mức độ tin cậy rằng $x$ là \textit{bonafide}. Quyết định cuối cùng được đưa ra bằng cách so sánh score với một ngưỡng $\tau$: nếu $s \geq \tau$ thì phân loại là \textit{bonafide}, ngược lại là \textit{spoof}. Trong toàn bộ báo cáo này, nhãn được quy ước là \textit{bonafide} = 1 và \textit{spoof} = 0 để thống nhất với cách tính EER ở Chương \ref{chap:experiment}.

Trong bối cảnh Automatic Speaker Verification (ASV), CM thường được đặt trước hoặc song song với hệ thống xác thực người nói. ASV trả lời câu hỏi ``người nói có đúng là danh tính đã khai báo không?'', còn CM trả lời câu hỏi ``tín hiệu đầu vào có phải là giọng thật không?''. Hai bài toán liên quan nhưng không đồng nhất: một audio giả có thể bắt chước đúng speaker target và qua được ASV, nhưng vẫn cần bị CM phát hiện là \textit{spoof}. Vì vậy các challenge ASVspoof đánh giá CM như một thành phần bảo vệ ASV trước các tấn công tổng hợp, chuyển đổi giọng hoặc replay \cite{asvspoof2019,asvspoof2021}.

Một hệ thống CM điển hình gồm ba bước. \textbf{Front-end} chuyển waveform thành representation: có thể là hand-crafted spectral features như LFCC/CQCC, feature học trực tiếp từ raw waveform, hoặc representation từ self-supervised learning (SSL) model như Wav2Vec2, XLS-R, WavLM. \textbf{Back-end} nhận representation này để mô hình hoá quan hệ thời gian, phổ hoặc speaker/utterance-level cue. \textbf{Classifier/scoring head} sinh score cuối cùng cho từng utterance. Cách tách front-end và back-end này giúp so sánh các nhóm kiến trúc ở Phần \ref{sec:architectures}: khác biệt chính giữa LFCC+LCNN, AASIST và XLS-R+Nes2Net không chỉ nằm ở classifier, mà nằm ở loại representation mà mô hình được phép khai thác.

Về threat model, các tấn công trong SDD thường được chia thành bốn nhóm chính:
\begin{itemize}
    \item \textbf{Logical Access (LA).} Attacker tiêm audio giả trực tiếp vào pipeline số, không cần phát lại qua loa hoặc microphone. Hai dạng phổ biến là TTS và VC. TTS sinh giọng nói từ văn bản; VC giữ nội dung hoặc đặc tính ngữ âm của nguồn nhưng biến đổi speaker identity sang giọng đích. LA là kịch bản chính của ASVspoof 2019 LA, ASVspoof 2021 DF và ASVspoof 5 Track 1.
    \item \textbf{Physical Access (PA) / replay.} Attacker phát lại audio qua thiết bị vật lý rồi thu hoặc đưa vào hệ thống qua microphone. Khác với LA, replay chứa thêm dấu vết của loa, phòng, microphone và khoảng cách thu. Do đó một detector huấn luyện cho LA không nhất thiết generalize sang PA, và ngược lại.
    \item \textbf{Deepfake/in-the-wild setting.} Audio được thu thập từ social media, video streaming hoặc nguồn công khai, thường đã qua codec, noise, trimming, post-processing và có metadata không đầy đủ. Đây không phải một attack mechanism đơn lẻ, mà là một deployment-like condition dùng để kiểm tra cross-domain generalization.
    \item \textbf{Adversarial và neural codec attacks.} Adversarial attacks thêm perturbation có chủ đích để làm sai lệch detector. Neural codec attacks liên quan đến các TTS/VC pipeline hiện đại sử dụng codec representation như EnCodec/SNAC/Vocos, khiến artifact không còn giống các spectral artifact truyền thống trong benchmark cũ. ASVspoof 5 đưa các yếu tố này vào benchmark ở quy mô lớn hơn các phiên bản trước \cite{asvspoof5_arxiv2024}.
\end{itemize}

Báo cáo này tập trung thực nghiệm vào nhóm LA và deepfake/in-the-wild vì bốn dataset reproduce ở Chương \ref{chap:experiment} chủ yếu thuộc hai nhóm này. PA/replay và adversarial robustness được trình bày ở mức cơ sở lý thuyết để làm rõ bức tranh nghiên cứu, nhưng chưa phải trọng tâm thực nghiệm của Thực tập 1.

\section{Các tập dữ liệu benchmark}

Dataset trong SDD không chỉ khác nhau về số lượng sample, mà còn khác nhau về giả định đánh giá. Có thể phân loại theo bốn trục chính: attack scenario, recording/deployment condition, attack generation family và language coverage. Bốn dataset được reproduce trong báo cáo không nhằm bao trùm toàn bộ landscape, mà đại diện cho bốn stress test chính trong scope Thực tập 1: clean benchmark, codec robustness, modern ASVspoof attacks và in-the-wild domain shift.

Bảng \ref{tab:sdd-datasets} tổng hợp các dataset chính. Cột ``Kích thước'' ưu tiên ghi split hoặc subset được dùng trong báo cáo nếu dataset được reproduce; với dataset chỉ tham chiếu, giá trị được ghi theo paper gốc ở mức tổng quan. Các dataset không reproduce được giữ trong bảng khi chúng đại diện cho một threat model quan trọng nhưng lệch scope hoặc chưa cần thiết cho pipeline thí nghiệm hiện tại.

\begin{table}[H]
    \centering
    \footnotesize
    \begin{adjustbox}{width=\textwidth}
    \begin{tabular}{p{2.4cm}|p{1.5cm}|p{2.2cm}|p{1.5cm}|p{3.4cm}|p{3.8cm}|p{4.2cm}}
        \textbf{Dataset} & \textbf{Năm} & \textbf{Kịch bản} & \textbf{Ngôn ngữ} & \textbf{Kích thước dùng/tham chiếu} & \textbf{Đặc điểm chính} & \textbf{Vai trò trong báo cáo}
        \\ \hline
        ASVspoof 2019 LA & 2019 & LA, TTS/VC & Anh & 71,237 utterances trong eval split dùng ở project & Clean lab, 19 attacks, metadata attack rõ & Reproduce: benchmark baseline trong điều kiện sạch
        \\ \hline
        ASVspoof 2021 DF & 2021 & Deepfake/LA + codec & Anh & 458,868 utterances trong eval subset dùng ở project & Audio kế thừa ASV2019, re-encode qua nhiều codec/bitrate & Reproduce: stress test cho codec robustness
        \\ \hline
        ASVspoof 5 Track 1 & 2024/2025 & LA/deepfake + adversarial context & Đa ngôn ngữ & 140,950 utterances trong dev split dùng ở project & Crowdsourced speech, 32 attack algorithms, surrogate/adversarial setting & Reproduce: trọng tâm preliminary error analysis; cần bổ sung full eval ở bước sau \cite{asvspoof5_design,asvspoof5_eval}
        \\ \hline
        In-the-Wild & 2022 & In-the-wild deepfake & Chủ yếu Anh & 31,779 utterances, 58 speakers trong bản dùng ở project & Scraped từ nguồn công khai, 20.8h bonafide và 17.2h spoofed audio & Reproduce: probe cho cross-domain generalization \cite{in_the_wild}
        \\ \hline
        MLAAD v9 & 2024/2026 & Multi-lingual TTS & 51 ngôn ngữ & 678.3h synthetic voice theo arXiv v9 & 140 TTS models, 78 architectures & Không reproduce: phù hợp cross-lingual extension, nhưng chưa phải câu hỏi chính của Thực tập 1 \cite{mlaad}
        \\ \hline
        CodecFake / CodecFake+ & 2024/2025 & Codec-based speech generation & Nhiều nguồn & CodecFake+ dùng 31 codec models cho train và 17 CoSG models cho eval & Nhắm trực tiếp vào neural codec / CoSG attacks & Không reproduce: phù hợp để follow-up neural codec robustness, nhưng ASVspoof 5 đang là benchmark hiện đại chính \cite{codecfake,codecfakeplus}
        \\ \hline
        EchoFake & 2025 & Zero-shot TTS + physical replay & Anh & Hơn 120h audio, hơn 13K speakers & Replay-aware practical SDD với device/environment đa dạng & Không reproduce: quan trọng cho deployment/replay, nhưng lệch LA/DF scope hiện tại \cite{echofake}
        \\ \hline
        ADD 2022/2023 & 2022/2023 & Low-quality, partial fake, fake game, localization & Chủ yếu Mandarin & Nhiều track challenge & Mở rộng ngoài binary utterance-level detection & Không reproduce: task/protocol khác với pipeline EER utterance-level \cite{add2022,add2023}
        \\ \hline
        PartialSpoof & 2021/2022 & Partially spoofed utterances & Anh & Dựa trên ASVspoof 2019 LA & Có utterance-level và segment-level labels & Không reproduce: cần localization/segment-level setup, không khớp thí nghiệm hiện tại \cite{partialspoof}
        \\ \hline
        WaveFake & 2021 & TTS/vocoder & Anh + Nhật & Khoảng 104K utterances theo paper & Tập trung GAN/neural vocoder artifacts & Tham chiếu lịch sử cho vocoder artifacts; không đủ hiện đại làm main benchmark
        \\ \hline
        FoR & 2019 & TTS & Anh & Khoảng 198K utterances theo paper & Dataset real/fake speech đời đầu, nhiều TTS cũ & Tham chiếu lịch sử, không dùng làm benchmark chính
    \end{tabular}
    \end{adjustbox}
    \caption{Tổng hợp một số dataset SDD tiêu biểu.}
    \label{tab:sdd-datasets}
\end{table}

\textbf{ASVspoof 2019 LA} là điểm xuất phát hợp lý vì điều kiện sạch, protocol rõ và được dùng rộng rãi trong cộng đồng ASV anti-spoofing \cite{asvspoof2019}. Điểm mạnh của nó là metadata attack type đầy đủ, cho phép phân tích từng nhóm TTS/VC. Hạn chế là nhiều attack phản ánh công nghệ synthesis trước 2019 và không có codec/post-processing phức tạp. Vì vậy EER thấp trên ASVspoof 2019 LA không đủ để kết luận mô hình sẽ hoạt động tốt ngoài môi trường benchmark.

\textbf{ASVspoof 2021 DF} được thiết kế để đưa yếu tố codec và transmission condition vào đánh giá \cite{asvspoof2021}. Về bản chất, dataset này dùng lại nguồn tấn công từ ASVspoof 2019 nhưng xử lý qua nhiều codec/bitrate. Thiết kế này hữu ích vì cô lập được một câu hỏi cụ thể: khi spectral artifacts bị nén hoặc làm mờ, detector còn phân biệt được \textit{bonafide}/\textit{spoof} không? Đây là lý do ASVspoof 2021 DF được dùng trong báo cáo như một stress test cho codec robustness.

\textbf{ASVspoof 5} mở rộng benchmark sang crowdsourced speech, nhiều speaker/recording condition hơn, 32 attack algorithms và adversarial attacks được đưa vào lần đầu trong chuỗi ASVspoof \cite{asvspoof5_arxiv2024,asvspoof5_design}. Paper đánh giá challenge năm 2026 cho thấy nhiều hệ thống vẫn suy giảm dưới adversarial attacks và neural encoding/compression, nên ASVspoof 5 là nguồn evidence quan trọng cho modern attack robustness \cite{asvspoof5_eval}. Trong phạm vi Thực tập 1, báo cáo dùng Track 1 dev split vì split này có metadata phù hợp cho phân tích lỗi sơ bộ. Việc chưa hoàn thành full eval split được ghi rõ ở Phần \ref{sec:error-analysis}; do đó các kết luận từ ASVspoof 5 trong báo cáo cần được hiểu là preliminary diagnosis, không phải đánh giá cuối cùng.

\textbf{In-the-Wild} khác với chuỗi ASVspoof ở chỗ nó không cố kiểm soát attack generator. Dataset được thu thập từ nguồn công khai cho 58 public figures, gồm 20.8 giờ \textit{bonafide} và 17.2 giờ spoofed audio \cite{in_the_wild}. Điểm mạnh của nó là phản ánh real-world condition: codec không đồng nhất, chất lượng nguồn khác nhau, speaker distribution khác training data và metadata attack không đầy đủ. Điểm yếu cũng nằm ở chính điều này: khó phân tích lỗi theo attack type và có khả năng tồn tại label noise. Trong báo cáo, In-the-Wild được dùng như probe cho domain shift, không dùng để rút ra kết luận chi tiết về từng attack mechanism.

Các dataset ngoài nhóm reproduce giúp làm rõ phần landscape còn lại. MLAAD v9 đáng chú ý cho cross-lingual evaluation vì mở rộng lên 51 ngôn ngữ và 140 TTS models \cite{mlaad}. CodecFake/CodecFake+ nhắm trực tiếp vào codec-based speech generation, một hướng tấn công mới mà nhiều detector train trên vocoder-era datasets có thể bỏ sót \cite{codecfake,codecfakeplus}. EchoFake đưa yếu tố physical replay vào speech deepfake thực tế \cite{echofake}. ADD và PartialSpoof mở rộng bài toán sang low-quality, partial và localized manipulation \cite{add2022,add2023,partialspoof}. Những dataset này chưa được đưa vào reproduce vì lệch task, đòi hỏi protocol/metric khác hoặc vượt tài nguyên của Thực tập 1; tuy nhiên chúng nên được giữ trong literature review và future work để tránh hiểu nhầm rằng bốn dataset reproduce là đầy đủ cho mọi kịch bản SDD.

\section{Các nhóm kiến trúc} \label{sec:architectures}

Các kiến trúc SDD có thể được nhìn theo câu hỏi: mô hình lấy thông tin gì từ audio, representation đó đến từ đâu và hệ thống có cơ chế nào để giữ ổn định khi attack/domain thay đổi hay không. Ở tầng cao, có thể gom thành bốn nhóm lớn thay vì liệt kê ngay từng model.

\textbf{Signal/task-specific supervised detectors.} Nhóm này dùng representation được thiết kế hoặc chuẩn hoá theo signal processing, rồi huấn luyện classifier/back-end supervised trên dữ liệu anti-spoofing. Nhánh cổ điển gồm LFCC, CQCC, MFCC hoặc IMFCC với GMM/SVM/LCNN. Nhánh học sâu hơn dùng spectrogram, CQT hoặc LFCC làm input cho LCNN, ResNet, TDNN, Transformer/Conformer. Ưu điểm là nhẹ, dễ triển khai và có baseline lâu đời trong ASVspoof. Điểm yếu là representation thường nhạy với artifact của training dataset: khi audio bị codec, replay channel hoặc neural codec generation làm thay đổi spectral/phase cues, khả năng generalize có thể giảm. LFCC+LCNN trong báo cáo đại diện cho nhóm này.

\textbf{Raw waveform / end-to-end detectors.} RawNet2 và AASIST bỏ qua feature engineering thủ công, học trực tiếp từ waveform \cite{rawnet2,aasist}. RawNet2 dùng SincConv/residual blocks để học filter và temporal pattern. AASIST bổ sung heterogeneous graph attention để mô hình hoá quan hệ giữa spectral branch và temporal branch; ý tưởng chính là artifact của spoofed speech có thể xuất hiện trong quan hệ spectro-temporal dài hơn, không chỉ ở một frame hoặc một vùng tần số riêng lẻ. AASIST/AASIST-L là mốc quan trọng vì đạt kết quả tốt trên ASVspoof 2019 LA với số tham số nhỏ \cite{aasist}. Tuy nhiên, vì không có large-scale speech pre-training, nhóm này vẫn có nguy cơ học artifact gắn với training distribution.

\textbf{SSL/foundation front-end + anti-spoofing back-end.} Nhóm này dùng Wav2Vec2, XLS-R, WavLM, HuBERT hoặc speech encoder lớn làm front-end, sau đó gắn back-end như AASIST, MFA, Nes2Net, pooling/MLP hoặc layer-selection module \cite{wav2vec2,xlsr,wavlm}. Lợi thế kỳ vọng là SSL representation chứa prior về phonetic structure, speaker/channel variation và acoustic regularity từ pre-training quy mô lớn, nên có nền tảng ổn định hơn khi attack hoặc codec thay đổi. Các paper tiêu biểu gồm Wav2Vec2/XLS-R + AASIST của Tak và cộng sự \cite{ssl_aasist}, WavLM + Multi-Fusion Attentive classifier của Guo và cộng sự \cite{wavlm_mfa}, AASIST3 \cite{aasist3}, WavLM back-ends/fusion trong ASVspoof 5 \cite{wavlm_backend,wavlm_ensemble} và Nes2Net \cite{nes2net}. Hạn chế chính là chi phí: riêng XLS-R 300M hoặc WavLM Large đã lớn hơn nhiều so với AASIST/LFCC+LCNN, và hiệu năng phụ thuộc checkpoint, layer selection, fine-tuning, augmentation và calibration.

\textbf{System-level robustness strategies.} Các hệ thống mạnh gần đây thường không chỉ thay backbone mà còn dùng codec/noise/reverb augmentation, CodecFake-aware training, score calibration và score-level/feature-level fusion. ASVspoof 5 evaluation 2026 cho thấy nhiều hệ thống vẫn suy giảm dưới adversarial attacks và neural encoding/compression \cite{asvspoof5_eval}, nên robustness cần được xem như một tầng thiết kế hệ thống. Whisper-based detector hoặc Whisper+AASIST là hướng emerging cho foundation-model front-end \cite{whisper_aasist}, nhưng trong literature hiện tại Wav2Vec2/XLS-R/WavLM vẫn là các encoder được dùng rộng rãi hơn trong anti-spoofing benchmark.

\begin{table}[H]
    \centering
    \footnotesize
    \begin{adjustbox}{width=\textwidth}
    \begin{tabular}{p{3.0cm}|p{4.0cm}|p{5.0cm}|p{4.0cm}}
        \textbf{Nhóm} & \textbf{Paper/dataset đại diện} & \textbf{Vai trò trong review} & \textbf{Trạng thái trong báo cáo}
        \\ \hline
        Modern benchmark & ASVspoof 5 design/evaluation \cite{asvspoof5_design,asvspoof5_eval} & Crowdsourced speech, adversarial setting, neural encoding/compression & Reproduce một phần qua Track 1 dev; full eval còn pending
        \\ \hline
        Codec-based attacks & CodecFake, CodecFake+ \cite{codecfake,codecfakeplus} & Chỉ ra gap với CoSG/neural codec speech & Chưa reproduce; đưa vào future work
        \\ \hline
        Practical replay & EchoFake \cite{echofake} & Zero-shot TTS + physical replay dưới device/environment đa dạng & Chưa reproduce; lệch LA/DF scope
        \\ \hline
        Partial/localized fake & ADD 2022/2023, PartialSpoof \cite{add2022,add2023,partialspoof} & Mở rộng sang low-quality, partial fake, localization & Chưa reproduce; cần task/metric khác
        \\ \hline
        SSL back-end design & WavLM+MFA, WavLM back-ends/ensemble \cite{wavlm_mfa,wavlm_backend,wavlm_ensemble} & Đại diện cho hướng SSL + fusion gần đây & Chưa reproduce; nên cân nhắc sau XLS-R baseline
        \\ \hline
        Reproduced SSL systems & AASIST3, XLS-R+AASIST, Nes2Net \cite{aasist3,ssl_aasist,nes2net} & Liên hệ trực tiếp với sáu mô hình Chương \ref{chap:experiment} & Đã đưa vào bảng EER
    \end{tabular}
    \end{adjustbox}
    \caption{Một số nguồn đại diện cần nhắc trong literature review.}
    \label{tab:literature-sources}
\end{table}

Bảng \ref{tab:reproduced-models} tổng hợp sáu mô hình được reproduce trong Chương \ref{chap:experiment}. Các giá trị hiệu năng trong cột ``Kết quả paper'' chỉ để đặt mô hình vào bối cảnh, không dùng thay cho kết quả reproduce.

\begin{table}[H]
    \centering
    \footnotesize
    \begin{adjustbox}{width=\textwidth}
    \begin{tabular}{p{2.5cm}|p{2.7cm}|p{2.6cm}|p{2.1cm}|p{3.7cm}|p{3.7cm}|p{3.3cm}}
        \textbf{Mô hình} & \textbf{Input/front-end} & \textbf{Back-end} & \textbf{Compute cost} & \textbf{Điểm mạnh kỳ vọng} & \textbf{Failure mode cần chú ý} & \textbf{Kết quả paper}
        \\ \hline
        LFCC + LCNN & LFCC spectral features & LCNN & Thấp & Nhẹ, dễ train, phù hợp baseline & Nhạy với codec và unseen artifacts & LFCC/LCNN variants thường là baseline trong ASVspoof
        \\ \hline
        AASIST & Raw waveform & Spectro-temporal graph attention & Thấp-trung bình & Compact, học cue trực tiếp từ waveform & Có thể overfit vào artifact của ASV19-style attacks & 0.83\% EER trên ASVspoof 2019 LA \cite{aasist}
        \\ \hline
        AASIST-L & Raw waveform & Lightweight AASIST & Thấp & Giảm tham số, phù hợp constrained setting & Capacity thấp hơn AASIST, vẫn thiếu SSL prior & 0.99\% EER trên ASVspoof 2019 LA \cite{aasist}
        \\ \hline
        AASIST3 & Wav2Vec2 SSL features & KAN bridge + AASIST & Cao & Khai thác SSL representation và regularization & Phụ thuộc checkpoint/condition ASVspoof 2024; nặng & Báo cáo cho ASVspoof 2024 \cite{aasist3}
        \\ \hline
        XLS-R + AASIST & XLS-R 300M & AASIST & Cao & Cross-lingual SSL prior, back-end đã quen với spoofing & Tốn VRAM; hiệu năng phụ thuộc fine-tuning/checkpoint & Kết quả mạnh trên ASVspoof 2021 LA/DF \cite{ssl_aasist}
        \\ \hline
        XLS-R + Nes2Net-X & XLS-R 300M & Nes2Net-X & Cao, back-end nhẹ & Tận dụng multi-layer SSL features với back-end gọn & Front-end vẫn chiếm chi phí chính & Kết quả tốt trên ASVspoof 2021 và In-the-Wild \cite{nes2net}
    \end{tabular}
    \end{adjustbox}
    \caption{So sánh sáu mô hình được tái lập.}
    \label{tab:reproduced-models}
\end{table}

\section{Phương pháp đánh giá}

Metric chính của báo cáo là EER. Với một ngưỡng $\tau$, mô hình phân loại các utterance có score $s \geq \tau$ là \textit{bonafide} và $s < \tau$ là \textit{spoof}. Khi đó:
\begin{equation}
    \mathrm{FAR}(\tau) = \frac{\mathrm{FP}(\tau)}{\mathrm{FP}(\tau) + \mathrm{TN}(\tau)}
\end{equation}
\begin{equation}
    \mathrm{FRR}(\tau) = \frac{\mathrm{FN}(\tau)}{\mathrm{FN}(\tau) + \mathrm{TP}(\tau)}
\end{equation}

Ở đây FAR (False Acceptance Rate) là tỉ lệ \textit{spoof} bị chấp nhận nhầm như \textit{bonafide}, còn FRR (False Rejection Rate) là tỉ lệ \textit{bonafide} bị từ chối nhầm như \textit{spoof}. EER là giá trị tại ngưỡng $\tau^*$ sao cho:
\begin{equation}
    \mathrm{FAR}(\tau^*) \approx \mathrm{FRR}(\tau^*).
\end{equation}
Trong thực tế, $\tau^*$ được tìm bằng cách sweep qua các score hoặc nội suy trên ROC/DET curve.

Ưu điểm của EER là không phụ thuộc vào một operating threshold cố định, nên phù hợp để so sánh các mô hình có thang score khác nhau. Đây là lý do EER được dùng làm metric chính trong Chương \ref{chap:experiment}, nơi mục tiêu là reproduce và so sánh tương đối giữa sáu mô hình trên bốn dataset. Tuy nhiên EER cũng có hạn chế: nó giả định hai loại lỗi có chi phí cân bằng, trong khi deployment thực tế có thể xem spoof false acceptance nguy hiểm hơn bonafide false rejection. EER cũng không phản ánh calibration của score; một mô hình có EER thấp vẫn có thể tạo score không tương ứng với xác suất thật, gây khó khi chọn threshold cố định.

Các challenge ASVspoof còn dùng minimum tandem Detection Cost Function (min t-DCF) để đánh giá CM khi đặt trong hệ thống ASV hoàn chỉnh \cite{tdcf}. Metric này tính đến prior và cost của nhiều loại lỗi trong pipeline ASV+CM, nên gần deployment hơn EER trong bối cảnh speaker verification. Tuy nhiên min t-DCF cần thông tin về ASV system, ASV score và cost model. Trong các kết quả reproduce của báo cáo, dữ liệu hiện có chỉ gồm CM scores và labels, vì vậy EER là lựa chọn nhất quán hơn cho phạm vi Thực tập 1.

Ngoài EER, ROC curve và DET curve là công cụ trực quan để xem trade-off giữa các loại lỗi ở nhiều threshold. ROC biểu diễn True Positive Rate theo False Positive Rate; DET biểu diễn FAR theo FRR trên thang xác suất Gaussian, thường được cộng đồng speaker verification và anti-spoofing dùng vì dễ quan sát vùng lỗi thấp. Calibration metrics như log loss hoặc Expected Calibration Error có ý nghĩa khi hệ thống cần score xác suất hoặc threshold cố định, nhưng chưa được phân tích sâu trong báo cáo này.

\section{Các thách thức hiện tại}

Từ các survey gần đây và từ thiết kế của các benchmark ASVspoof/In-the-Wild, có thể rút ra bốn thách thức chính làm nền cho phần thực nghiệm ở Chương \ref{chap:experiment} \cite{asvspoof5_arxiv2024,in_the_wild,sdd_survey}.

\textbf{Cross-domain generalization.} Dataset evidence rõ nhất là khoảng cách giữa clean benchmark và in-the-wild benchmark. ASVspoof 2019 LA có điều kiện kiểm soát, attack metadata rõ và ít yếu tố channel phức tạp. In-the-Wild lại chứa audio từ nguồn công khai, speaker distribution khác, codec khác và attack metadata thiếu. Nếu mô hình học artifact phụ thuộc dataset thay vì cue tổng quát của spoofed speech, EER sẽ tăng mạnh khi chuyển domain. Về mặt kiến trúc, hand-crafted features và small end-to-end models dễ bị ảnh hưởng hơn vì representation được học trong phạm vi dữ liệu anti-spoofing hẹp. SSL front-end được kỳ vọng giảm một phần gap này nhờ pre-training trên audio đa dạng, nhưng cần kết quả cross-dataset ở Chương \ref{chap:experiment} để kiểm chứng trong project cụ thể.

\textbf{Codec robustness.} ASVspoof 2021 DF được thiết kế trực tiếp cho vấn đề này: audio từ nguồn ASVspoof 2019 được xử lý qua nhiều codec/bitrate \cite{asvspoof2021}. Codec lossy có thể xoá high-frequency components, làm mờ phase/spectral artifacts hoặc tạo artifact mới không liên quan đến synthesis. Điều này ảnh hưởng đặc biệt tới các mô hình dựa vào spectral cues cục bộ như LFCC+LCNN, và cũng có thể ảnh hưởng tới raw waveform models nếu chúng học artifact ở waveform distribution gốc. Nếu SSL representation học được cấu trúc speech ổn định hơn, các mô hình XLS-R-based trong Chương \ref{chap:experiment} nên có generalization gap nhỏ hơn trên ASVspoof 2021 DF.

\textbf{Modern attacks và adversarial setting.} Các hệ TTS/VC hiện đại dùng neural codec, diffusion hoặc pipeline nhiều tầng có artifact khác với các hệ cũ trong ASVspoof 2019. ASVspoof 5 đưa thêm các attack gần thời điểm hiện tại hơn và có adversarial track riêng \cite{asvspoof5_arxiv2024,asvspoof5_design,asvspoof5_eval}. CodecFake/CodecFake+ cũng cho thấy deepfake từ codec-based speech generation có thể tạo ra failure mode khác với vocoder-era datasets \cite{codecfake,codecfakeplus}. Thách thức ở đây không chỉ là thêm nhiều attack, mà là thay đổi bản chất của cue: detector có thể không còn tìm thấy dấu hiệu vocoder/spectral quen thuộc. Vì vậy thứ hạng mô hình trên ASVspoof 2019 LA không nhất thiết dự đoán thứ hạng trên ASVspoof 5 hoặc CodecFake-style data.

\textbf{Multi-linguality, fairness và deployment constraints.} Phần lớn benchmark kinh điển tập trung tiếng Anh hoặc một số corpus giới hạn. MLAAD v9 mở rộng câu hỏi sang 51 ngôn ngữ, cho thấy SDD không nên chỉ được đánh giá trên English-centric data \cite{mlaad}. EchoFake lại nhấn mạnh ràng buộc deployment khác: audio synthetic có thể đi qua replay channel và thiết bị tiêu dùng trước khi đến detector \cite{echofake}. Ngoài ra, deployment thực tế còn bị ràng buộc bởi latency, VRAM, privacy và khả năng chạy trên edge device. AASIST/AASIST-L có lợi thế compute, nhưng có thể kém ổn định khi domain thay đổi; XLS-R/Nes2Net có lợi thế representation nhưng chi phí cao. Trade-off này không được giải quyết chỉ bằng một bảng EER, mà cần phân tích đồng thời accuracy, robustness và resource cost trong các bước tiếp theo.

Các thách thức trên tạo khung diễn giải cho Chương \ref{chap:experiment}: ASVspoof 2019 LA kiểm tra clean benchmark performance, ASVspoof 2021 DF kiểm tra codec robustness, ASVspoof 5 kiểm tra modern attack robustness ở mức sơ bộ, và In-the-Wild kiểm tra cross-domain generalization.

\chapter{Reproduce và phân tích kết quả} \label{chap:experiment}

\section{Thiết lập thí nghiệm}

Toàn bộ thí nghiệm trong báo cáo được thực hiện trên nền tảng Kaggle với GPU P100 hoặc T4. Việc lựa chọn Kaggle làm môi trường tính toán là do nền tảng cung cấp GPU miễn phí phù hợp với quy mô tái lập, đặc biệt là suy luận mô hình SSL 300M trên các tập eval cỡ vài chục đến vài trăm nghìn utterance. Bên cạnh đó, Kaggle cho phép quản lý dataset và checkpoint dưới dạng Kaggle Datasets công khai, hỗ trợ tính khả lập.

\textbf{Mã nguồn và cấu trúc.} Mỗi dataset được đánh giá bằng một notebook riêng: \texttt{notebooks/eval\_asvspoof\_2019.ipynb}, \texttt{notebooks/eval\_asvspoof\_2021.ipynb}, \texttt{notebooks/eval\_asvspoof\_5.ipynb} và \texttt{notebooks/eval\_in\_the\_wild.ipynb}. Các notebook này được sinh tự động từ một template chung trong \texttt{scripts/create\_eval\_notebooks.py}. Cách tổ chức này giúp giữ logic suy luận nhất quán giữa các dataset, bao gồm cách load checkpoint và cách tính EER, trong khi vẫn cho phép tuỳ biến phần parser metadata cho từng dataset.

\textbf{Nguồn checkpoint.} Tất cả các mô hình được sử dụng ở chế độ pretrained, không thực hiện fine-tuning thêm trong phạm vi báo cáo này. Các checkpoint được lấy từ các nguồn công khai như sau: AASIST và AASIST-L từ repository chính thức của tác giả; AASIST3 từ HuggingFace \texttt{MTUCI/AASIST3}; XLS-R 300M từ release Fairseq; XLS-R + AASIST từ checkpoint công bố bởi nhóm tác giả SSL-AASIST; XLS-R + Nes2Net từ release của nhóm tác giả Nes2Net; LFCC+LCNN được huấn luyện lại từ đầu trên tập huấn luyện ASVspoof 2019 LA khoảng 10 epoch do không có checkpoint công khai phù hợp.

\textbf{Cấu trúc kết quả.} Toàn bộ điểm số đầu ra được lưu vào các file \texttt{results/<dataset>/results.pkl}, mỗi file là một \texttt{dict} Python với cấu trúc:

\begin{verbatim}
{
    "<model_name>": {
        "eer":    float,       # EER (%)
        "scores": np.ndarray,  # (N,) - diem bonafide
        "labels": np.ndarray,  # (N,) - 1=bonafide, 0=spoof
    },
    ...
}
\end{verbatim}

Cấu trúc đồng nhất này cho phép các bước phân tích hậu kỳ như error analysis, score correlation và cross-dataset comparison hoạt động trên mọi dataset mà không cần xử lý đặc thù.

\textbf{Hướng mở rộng.} Để bổ sung một mô hình mới, cần thêm hàm load checkpoint và một cell suy luận tương ứng vào \texttt{scripts/create\_eval\_notebooks.py} rồi sinh lại các notebook. Để bổ sung một dataset mới, cần thêm parser protocol và đường dẫn audio vào cùng file generator. Cách tổ chức này được thiết kế để giảm chi phí mở rộng ở Internship 2.

\section{Kết quả EER tổng hợp}

Bảng \ref{tab:eer-results} tổng hợp EER (\%) của sáu mô hình trên bốn dataset. Các giá trị được tính trên tập eval của từng dataset bằng \texttt{sklearn} từ điểm số \textit{bonafide} và nhãn ground truth, với quy ước 1 = \textit{bonafide} và 0 = \textit{spoof}.

\begin{table}[H]
    \centering
    \begin{tabular}{l|c|c|c|c}
        \textbf{Mô hình} & \textbf{ASV19 LA} & \textbf{ASV21 DF} & \textbf{ASV5 (Track 1)} & \textbf{In-the-Wild}
        \\ \hline
        LFCC + LCNN & 19.64 & 33.81 & 22.60 & 70.23
        \\ \hline
        AASIST & 4.59 & 17.70 & 37.81 & 41.79
        \\ \hline
        AASIST-L & 6.74 & 19.13 & 39.47 & 45.28
        \\ \hline
        AASIST3 & 20.83 & 29.18 & 19.03 & 40.12
        \\ \hline
        XLS-R + AASIST & 1.17 & -- & 2.55 & --
        \\ \hline
        XLS-R + Nes2Net & \textbf{0.45} & \textbf{2.93} & \textbf{1.80} & \textbf{5.57}
    \end{tabular}
    \caption{EER (\%) của sáu mô hình trên bốn dataset. Số nhỏ hơn là tốt hơn. Các ô ``--'' tương ứng với các đánh giá chưa hoàn thành tại thời điểm viết báo cáo.}
    \label{tab:eer-results}
\end{table}

Một số xu hướng có thể quan sát từ Bảng \ref{tab:eer-results}. Thứ nhất, dải EER trải rộng đáng kể giữa các mô hình trên cùng một dataset. Trên ASVspoof 2019 LA, khoảng cách giữa mô hình tốt nhất quan sát được là XLS-R + Nes2Net với 0.45\% và mô hình kém nhất là AASIST3 với 20.83\% lên tới khoảng 20 điểm phần trăm. Trên In-the-Wild, khoảng cách này còn rộng hơn khi XLS-R + Nes2Net đạt 5.57\% trong khi LFCC+LCNN đạt 70.23\%.

Thứ hai, EER có xu hướng tăng khi chuyển từ ASVspoof 2019 LA sang các dataset khác đối với phần lớn các mô hình, dù mức độ tăng khác nhau. Riêng AASIST3 và LFCC+LCNN không tuân theo xu hướng đơn điệu này, như sẽ được phân tích ở Phần \ref{sec:dataset-analysis}.

Thứ ba, hai mô hình dùng XLS-R 300M làm front-end, gồm XLS-R + AASIST và XLS-R + Nes2Net, duy trì EER ở dải thấp hơn đáng kể trên các dataset khó như ASVspoof 2021 DF, ASVspoof 5 và In-the-Wild so với các mô hình không dùng SSL front-end. Đây là quan sát chính của báo cáo và là dẫn dắt cho phần đề xuất ở Chương \ref{chap:conclusion}.

Cần lưu ý rằng các nhận xét trên chỉ giới hạn trong phạm vi sáu mô hình và bốn dataset được khảo sát. Việc kết luận tổng quát hơn, ví dụ về tính tổng quát hoá của SSL front-end nói chung, đòi hỏi phân tích lỗi chi tiết và mở rộng tập mô hình/dataset ở các giai đoạn tiếp theo.

\section{Phân tích EER theo từng dataset} \label{sec:dataset-analysis}

\textbf{ASVspoof 2019 LA.} Dataset này gồm 71,237 utterance, trong đó có 7,355 \textit{bonafide} và 63,882 \textit{spoof}. Đây là benchmark sạch nhất trong bốn dataset. Bốn mô hình đạt EER dưới 7\% gồm XLS-R + Nes2Net (0.45), XLS-R + AASIST (1.17), AASIST (4.59) và AASIST-L (6.74), trong khi LFCC+LCNN (19.64) và AASIST3 (20.83) cao hơn rõ rệt. Riêng AASIST3 cao bất ngờ trên dataset này; mô hình được công bố cho ASVspoof 2024 và có thể đã được huấn luyện theo phân bố tấn công khác với ASVspoof 2019 LA. Điều này cần được xác nhận thêm trong phần error analysis. LFCC+LCNN ở mức 19.64 phù hợp với baseline thấp đã biết của nhóm hand-crafted features khi không được fine-tune kỹ.

\textbf{ASVspoof 2021 DF.} Dataset này gồm 458,868 utterance, trong đó có 16,977 \textit{bonafide} và 441,891 \textit{spoof}. Dataset tái sử dụng nguồn audio của ASVspoof 2019 nhưng được re-encode qua nhiều cấu hình codec lossy. Các mô hình không dùng SSL có xu hướng tăng EER khi chuyển từ ASVspoof 2019 LA sang ASVspoof 2021 DF: AASIST từ 4.59 lên 17.70 (+13.1 điểm phần trăm), AASIST-L từ 6.74 lên 19.13 (+12.4 điểm phần trăm), LFCC+LCNN từ 19.64 lên 33.81 (+14.2 điểm phần trăm), AASIST3 từ 20.83 lên 29.18 (+8.3 điểm phần trăm). Trong khi đó, XLS-R + Nes2Net chỉ tăng từ 0.45 lên 2.93 (+2.5 điểm phần trăm). Quan sát này gợi ý rằng codec lossy có ảnh hưởng đáng kể tới các mô hình hand-crafted hoặc end-to-end nhỏ; tuy nhiên đây mới là quan sát ở mức tổng EER, cần error analysis theo loại codec để xác nhận giả thuyết.

\textbf{ASVspoof 5 Track 1.} Dataset này gồm 140,950 utterance, trong đó có 31,334 \textit{bonafide} và 109,616 \textit{spoof}. ASVspoof 5 chứa 32 loại tấn công bao gồm các hệ neural codec TTS và adversarial perturbation. Đáng chú ý, trên dataset này thứ tự các mô hình thay đổi rõ rệt so với ASVspoof 2019 LA: AASIST (37.81) và AASIST-L (39.47) cao hơn LFCC+LCNN (22.60) và AASIST3 (19.03). XLS-R + AASIST (2.55) và XLS-R + Nes2Net (1.80) tiếp tục giữ EER thấp. Sự đảo thứ tự này cho thấy hành vi của mô hình trên ASVspoof 5 phụ thuộc vào loại tấn công cụ thể và có thể không tỉ lệ thuận với hiệu năng trên ASVspoof 2019 LA. Đây là một trong những lý do ASVspoof 5 được chọn làm dataset trọng tâm cho phần error analysis ở Phần \ref{sec:error-analysis}.

\textbf{In-the-Wild.} Dataset này gồm 31,779 utterance, trong đó có 19,963 \textit{bonafide} và 11,816 \textit{spoof}. Đây là tập probe cho khả năng tổng quát hoá liên miền: không có mô hình nào trong báo cáo được huấn luyện trên dữ liệu mạng xã hội. Tất cả các mô hình đều cho EER cao hơn so với ASVspoof 2019 LA. LFCC+LCNN tăng mạnh nhất, từ 19.64 lên 70.23, trong khi XLS-R + Nes2Net giữ EER ở mức 5.57. Khoảng cách này gợi ý rằng các mô hình không dùng SSL front-end có thể đang học các đặc trưng phụ thuộc nhiều vào điều kiện huấn luyện. Tuy nhiên cũng cần lưu ý rằng In-the-Wild có ground truth đến từ xác minh thủ công và có thể chứa label noise, do đó EER cao có thể phản ánh một phần vấn đề chất lượng nhãn.

\section{Phân tích lỗi trên ASVspoof 5} \label{sec:error-analysis}

Phần này sẽ được hoàn thiện sau khi hoàn thành hai bước: đánh giá đầy đủ trên tập eval của ASVspoof 5, vì hiện báo cáo dùng tập dev với metadata phong phú hơn; và bổ sung mô hình Nes2Net không có XLS-R để cô lập đóng góp của SSL front-end. Khi đó, các tiểu mục dự kiến gồm:

\begin{itemize}
    \item \textbf{Phân phối điểm số}: biểu đồ phân phối điểm \textit{bonafide} theo nhãn cho từng mô hình, từ đó quan sát mức độ tách biệt giữa hai phân phối.
    \item \textbf{Lỗi theo loại tấn công}: EER nhóm theo \texttt{attack\_tag} của ASVspoof 5 để xác định các tấn công gây lỗi nhiều nhất cho mỗi mô hình.
    \item \textbf{Lỗi theo codec}: EER nhóm theo \texttt{codec}/\texttt{codec\_q} để định lượng ảnh hưởng của codec.
    \item \textbf{Confident errors}: các utterance bị phân loại sai với confidence cao; nếu có audio tương ứng, có thể đính kèm spectrogram đại diện.
    \item \textbf{Failure overlap giữa các mô hình}: so sánh các utterance bị tất cả mô hình phân loại sai với các utterance chỉ một mô hình sai, nhằm phân biệt universal hard examples và model-specific weakness.
\end{itemize}

Các quan sát từ Phần \ref{sec:error-analysis} sẽ là cơ sở cho phần đề xuất hướng tiếp cận ở Chương \ref{chap:conclusion}.

\section{Nhận xét liên dataset}

Bảng \ref{tab:generalization-gap} trình bày generalization gap -- chênh lệch EER giữa từng dataset khó hơn và benchmark cơ sở ASVspoof 2019 LA -- như một cách đo định lượng mức độ tổng quát hoá của từng mô hình.

\begin{table}[H]
    \centering
    \begin{tabular}{l|c|c|c}
        \textbf{Mô hình} & \textbf{ASV21 DF - ASV19} & \textbf{ASV5 - ASV19} & \textbf{ITW - ASV19}
        \\ \hline
        LFCC + LCNN & +14.17 & +2.96 & +50.59
        \\ \hline
        AASIST & +13.11 & +33.22 & +37.20
        \\ \hline
        AASIST-L & +12.39 & +32.73 & +38.54
        \\ \hline
        AASIST3 & +8.35 & -1.80 & +19.29
        \\ \hline
        XLS-R + AASIST & -- & +1.38 & --
        \\ \hline
        XLS-R + Nes2Net & \textbf{+2.48} & \textbf{+1.35} & \textbf{+5.12}
    \end{tabular}
    \caption{Generalization gap (delta EER, điểm phần trăm) so với ASVspoof 2019 LA. Số dương nghĩa là EER tăng khi chuyển sang dataset khó hơn.}
    \label{tab:generalization-gap}
\end{table}

Nhóm dùng XLS-R 300M làm front-end có generalization gap nhỏ nhất trong phạm vi các dataset được đo. Đặc biệt, XLS-R + Nes2Net giữ gap dưới 5.5 điểm phần trăm trên cả ba dataset khó. Điều này nhất quán với kỳ vọng rằng representation từ mô hình SSL lớn có thể bền hơn trước codec và domain shift so với feature thủ công hoặc mô hình raw waveform nhỏ.

AASIST và AASIST-L có gap lớn nhất trên ASVspoof 5, lần lượt là +33.22 và +32.73 điểm phần trăm. Quan sát này gợi ý rằng các mô hình này có thể đặc biệt nhạy cảm với loại tấn công trong ASVspoof 5; cần error analysis ở Phần \ref{sec:error-analysis} để xác nhận. Trong khi đó, AASIST3 có gap âm trên ASVspoof 5 (-1.80 điểm phần trăm), nghĩa là EER trên ASVspoof 5 thấp hơn trên ASVspoof 2019 LA. Quan sát này nhất quán với khả năng AASIST3 được huấn luyện theo phân bố gần ASVspoof 2024 hơn ASVspoof 2019.

LFCC+LCNN có gap lớn nhất trên In-the-Wild (+50.59 điểm phần trăm). Kết hợp với khả năng có label noise trong In-the-Wild đã đề cập ở Phần \ref{sec:dataset-analysis}, cần thận trọng khi diễn giải con số này. Các nhận xét trên giới hạn trong phạm vi sáu mô hình và bốn dataset được khảo sát, đồng thời phụ thuộc vào việc các checkpoint đều được sử dụng ở chế độ pretrained mà không fine-tune. Các quan sát chi tiết hơn về nguyên nhân đứng sau các generalization gap này sẽ được rút ra sau khi hoàn thành error analysis ở Phần \ref{sec:error-analysis}.

\chapter{Kết luận và hướng phát triển} \label{chap:conclusion}

\section{Tóm tắt kết quả}

Báo cáo Thực tập 1 đã thực hiện ba khối công việc: (i) hệ thống hoá literature speech deepfake detection theo bốn trục: bài toán và metric, dataset, kiến trúc và thách thức (Chương \ref{chap:background}); (ii) tái lập benchmark chính với năm mô hình đại diện (LFCC+LCNN, AASIST, AASIST-L, XLS-R+AASIST, XLS-R+Nes2Net) trên bốn dataset (ASVspoof 2019 LA, ASVspoof 2021 DF, ASVspoof 5 Track 1, In-the-Wild) qua một pipeline đánh giá thống nhất, đồng thời ghi nhận lý do loại checkpoint AASIST3 public khỏi phân tích chính; và (iii) phân tích lỗi diagnostic trên ASVspoof 5 để bóc tách thành phần tạo nên EER cao của các mô hình hiện đại.

Năm phát hiện chính đã hệ thống ở Phần 3.6 định hình hướng đề xuất:
\begin{itemize}
    \item \textbf{F1, F3 -- Codec là bottleneck chính trên ASV5}, không phải attack family. Hai mô hình XLS-R có EER trên subset \texttt{nocodec} chỉ 2.71--3.78\% nhưng tăng lên 35--42\% trên C04/C07; trong khi EER theo attack tag (AC1/AC2/AC3) gần đồng nhất.
    \item \textbf{F2 -- Codec coverage gap giữa ASV21 DF và ASV5.} Generalization gap của XLS-R+Nes2Net chỉ +2.48 pp trên ASV21 DF nhưng +21.13 pp trên ASV5; ASV21 DF không bao phủ đầy đủ codec/encoding conditions khó có trong ASV5. ASV21 DF vẫn là benchmark codec hợp lệ, nhưng không đủ một mình để đánh giá codec robustness deployment-grade.
    \item \textbf{F4 -- Failure overlap và score correlation thấp giữa các họ kiến trúc} gợi ý cơ hội fusion ở mức hệ thống.
    \item \textbf{F5 -- SSL front-end là điều kiện cần nhưng không đủ.} Hai mô hình cùng XLS-R 300M front-end chênh nhau đáng kể tuỳ dataset; robustness còn phụ thuộc back-end, fine-tuning recipe, checkpoint selection và calibration.
\end{itemize}

Ba phát hiện đầu (F1--F3) hội tụ về cùng một câu chuyện: ngay cả các mô hình mạnh nhất hiện nay vẫn thất bại có hệ thống ở các điều kiện codec không có trong huấn luyện, và bottleneck nằm ở training/representation distribution chứ không ở loại tấn công. Đây là motivation chính cho đề xuất ở Phần \ref{sec:proposal}.

\section{Đề xuất: Codec-aware fine-tuning với score-level consistency regularization} \label{sec:proposal}

\textbf{Tuyên bố đề xuất.} Báo cáo đề xuất một phương pháp huấn luyện codec-aware cho SSL-based speech deepfake detector, gồm hai thành phần kết hợp:
\begin{enumerate}
    \item \textbf{Codec-balanced data augmentation} ở mức waveform, dùng tập codec đa dạng và codec-balanced sampling.
    \item \textbf{Score-level codec consistency regularization}: auxiliary loss buộc decision không đổi giữa clean view và codec view của cùng một utterance.
\end{enumerate}

Backbone là XLS-R + Nes2Net-X, được chọn vì là baseline mạnh nhất tổng thể trên ASV19 LA, ASV21 DF và In-the-Wild trong Chương \ref{chap:experiment}. Đánh giá chính dựa trên một metric headline duy nhất: \textbf{codec robustness gap} trên ASV5.

\subsection{Động lực}

Kết quả Chương \ref{chap:experiment} cho thấy ba điểm hội tụ về codec robustness:
\begin{itemize}
    \item \textbf{Vấn đề cụ thể, định lượng được:} trên ASV5 Track 1 eval, XLS-R + Nes2Net và XLS-R + AASIST có EER tổng 21.58\% và 19.60\%, với codec robustness gap (định nghĩa ở Phần \ref{sec:eval-protocol}) khoảng 39 pp và 36 pp tương ứng. Phần lớn lỗi của các mô hình tốt nhất hiện tại nằm ở một vài codec.
    \item \textbf{Khoảng trống của benchmark hiện tại:} ASV21 DF không bao phủ đầy đủ codec khó có trong ASV5; mô hình tốt trên ASV21 DF vẫn có thể thất bại nặng trên ASV5 (F2). Việc đánh giá codec robustness deployment cần hơn một benchmark.
    \item \textbf{Bottleneck không chỉ nằm ở back-end:} EER gần đồng nhất theo attack tag (F3) cho thấy thay back-end đơn thuần (Nes2Net-X sang Nes2Net-LA hoặc XLSR-Mamba) có thể không đủ nếu không xử lý codec mismatch song song.
\end{itemize}

Đồng thời, literature ở Phần \ref{sec:architectures} cho thấy các kiến trúc 2024--2025 (Nes2Net-X, Nes2Net-LA, XLSR-Mamba) chủ yếu cải thiện thiết kế back-end trên benchmark sạch, ít công trình tập trung trực tiếp vào codec-aware training cho SSL-based detector. Codec augmentation đã xuất hiện rải rác trong các submission ASVspoof, nhưng (a) chi tiết thường không được công bố thống nhất, và (b) codec-consistency regularization cho SSL-based CM chưa được khảo sát có hệ thống. Đây là khoảng trống mà đề xuất nhắm đến.

\subsection{Phương pháp}

\textbf{Backbone.} XLS-R 300M (partial unfreeze ở các transformer layer cuối) + Nes2Net-X back-end, lấy từ release của Liu et al. (2025) làm điểm khởi đầu và fine-tune theo các điều kiện huấn luyện ở Phần \ref{sec:experiment-design}.

\textbf{Codec-balanced augmentation pipeline.} Với mỗi utterance huấn luyện $x$, sinh một phiên bản codec-augmented $x_c = \mathrm{codec}_k(x)$, trong đó $k$ được lấy mẫu từ tập codec tham chiếu $\mathcal{K}$ theo phân bố cân bằng (mỗi batch chứa hỗn hợp nhiều codec, không lệch về một codec đơn lẻ). Tập codec tham chiếu $\mathcal{K}$ bao gồm các codec phổ thông (MP3, AAC, Opus) ở nhiều bitrate, kèm các điều kiện narrowband và transcoding chain. Quan trọng: $\mathcal{K}$ được chia thành \textbf{train-codec set} $\mathcal{K}_{\text{tr}}$ và \textbf{held-out codec set} $\mathcal{K}_{\text{ho}}$ disjoint (xem Phần \ref{sec:eval-protocol}).

\textbf{Score-level codec consistency regularization (formulation chính).} Với cùng một utterance $x$, model sinh logit $z(x)$ và $z(x_c)$. Loss tổng:
\begin{equation}
    \mathcal{L}
    =
    \mathcal{L}_{\text{CM}}\big(z(x), y\big)
    +
    \mathcal{L}_{\text{CM}}\big(z(x_c), y\big)
    +
    \lambda \cdot \mathcal{L}_{\text{cons}}\big(z(x), z(x_c)\big).
\end{equation}
Trong đó $\mathcal{L}_{\text{CM}}$ là binary cross-entropy chuẩn, và $\mathcal{L}_{\text{cons}}$ là \textit{score-level consistency}: mặc định là MSE giữa hai logit, hoặc KL divergence giữa hai phân phối sigmoid. Hệ số $\lambda$ là hyper-parameter, được sweep trên dev split. Diễn giải của loss này là ràng buộc \textit{quyết định} không đổi sau codec, không ràng buộc embedding hay feature trung gian, nhằm tránh rủi ro xoá cue spoof do ép embedding-invariance quá mạnh.

\textbf{Variant -- supervised contrastive theo class (long-term).} Một biến thể chặt hơn về mặt formulation, dành cho future work: thay vì chỉ ràng buộc cặp $(x, x_c)$ cùng utterance, dùng supervised contrastive loss cross-utterance, kéo các view \textit{bonafide} (gồm cả codec-augmented) lại gần nhau, tách khỏi cụm \textit{spoof}, và ngược lại. Lưu ý: định nghĩa ``within-class consistency'' chỉ trên cặp $(x, x_c)$ cùng utterance sẽ trùng với $\beta$ do cặp luôn cùng nhãn; vì vậy variant đúng đòi hỏi mở rộng cross-utterance, phức tạp hơn và được dời xuống Phần \ref{sec:future-work}.

\textbf{Train data.} Phase chính dùng \textbf{ASVspoof 2019 LA train/dev} làm nguồn huấn luyện có kiểm soát; codec view được sinh on-the-fly hoặc pre-computed bằng augmentation pipeline ở trên. ASVspoof 2021 DF được giữ làm evaluation set cho codec robustness (vì challenge 2021 không phát hành DF training split công khai), trừ khi xác minh được một nguồn train-compatible riêng. Cả ba điều kiện huấn luyện B1/B2/B3 (Phần \ref{sec:experiment-design}) dùng cùng training source ASV19 LA, cùng số optimizer steps và cùng số utterance gốc, chỉ khác phần codec view và consistency loss. Lựa chọn này quan trọng để claim ``codec-aware training là nguồn cải thiện'' tách bạch khỏi các yếu tố nhiễu khác. ASV5 train/dev được dành cho phase 2 in-domain adaptation (optional).

\subsection{Thiết kế thí nghiệm} \label{sec:experiment-design}

Bảng \ref{tab:proposal-training-conditions} liệt kê các điều kiện huấn luyện. B0 là baseline không fine-tune, lấy thẳng từ Chương \ref{chap:experiment}. MUST gồm 3 run B1/B2/B3 trên cùng backbone XLS-R + Nes2Net-X. SHOULD lặp lại B1--B3 với XLS-R + AASIST để trả lời F5 (codec-aware có giúp cả back-end khác không). COULD là các mở rộng tuỳ thời gian, gồm cả supervised contrastive variant và sweep front-end/back-end.

\begin{table}[H]
    \centering
    \footnotesize
    \begin{adjustbox}{width=\textwidth}
    \begin{tabular}{p{2.6cm}|p{3.8cm}|p{3.8cm}|p{3.4cm}|p{4.2cm}}
        \textbf{Điều kiện} & \textbf{Backbone} & \textbf{Augmentation} & \textbf{Consistency loss} & \textbf{Phạm vi}
        \\ \hline
        B0 & XLS-R + Nes2Net-X & -- & -- & Pretrained, đã có ở Chương \ref{chap:experiment}
        \\ \hline
        B1 & XLS-R + Nes2Net-X & Không & Không & MUST -- fine-tune control
        \\ \hline
        B2 & XLS-R + Nes2Net-X & Codec-balanced ($\mathcal{K}_{\text{tr}}$) & Không & MUST -- augmentation only
        \\ \hline
        B3 & XLS-R + Nes2Net-X & Codec-balanced ($\mathcal{K}_{\text{tr}}$) & Score-level ($\beta$) & MUST -- đề xuất chính
        \\ \hline
        B1'/B2'/B3' & XLS-R + AASIST & Như B1/B2/B3 & Như B1/B2/B3 & SHOULD -- kiểm chứng F5
        \\ \hline
        C* & WavLM front-end, Nes2Net-LA, XLSR-Mamba, supervised contrastive variant & Varies & Varies & COULD -- extension
    \end{tabular}
    \end{adjustbox}
    \caption{Các điều kiện huấn luyện trong proposal.}
    \label{tab:proposal-training-conditions}
\end{table}

\textbf{So sánh kỳ vọng.} B1 đóng vai trò \textit{fine-tune control}. Không claim B1 vs B0 đo riêng ``đóng góp của fine-tune'' vì chuyển từ pretrained sang fine-tune đổi cả domain training và optimizer dynamics; B1 chỉ là baseline để các so sánh sau đứng được. B2 vs B1 đo đóng góp của codec augmentation; B3 vs B2 đo đóng góp của consistency regularization. Decomposition này tách rõ hai thành phần augmentation và consistency.

\textbf{Fairness budget.} Để các so sánh B1 vs B2 vs B3 sạch về ngân sách huấn luyện: cùng số optimizer steps, cùng số utterance base (cùng tập ASV19 LA train), cùng learning rate schedule và cùng số seed. B2 không được ``thắng'' B1 chỉ vì thấy thêm augmented samples; vì vậy số \textit{gradient updates} được giữ cố định, augmentation chỉ thay đổi nội dung của batch, không thay đổi số batch. B3 thêm consistency loss term nhưng không thêm utterance, để tách đóng góp của loss khỏi đóng góp của data.

\textbf{Compute và infrastructure.} Toàn bộ MUST (3 run) được thiết kế để chạy trên Kaggle T4 với XLS-R partial unfreeze, gradient checkpointing và mixed precision. SHOULD cần thêm 3 run; nếu vượt budget, SHOULD chuyển sang COULD.

\subsection{Evaluation protocol và metric headline} \label{sec:eval-protocol}

\textbf{Codec robustness gap (CodecGap)} là metric headline, định nghĩa trên ASV5 Track 1 eval:
\begin{equation}
    \mathrm{CodecGap}
    =
    \frac{\mathrm{EER}(C04) + \mathrm{EER}(C07)}{2}
    -
    \mathrm{EER}(\texttt{nocodec}).
\end{equation}
Trong đó EER theo nhóm được tính bằng chế độ \codepath{spoof_group_vs_all_bonafide} như Phần \ref{sec:error-analysis}. Lựa chọn mean(C04, C07) thay vì max ổn định hơn theo một codec đơn lẻ. Baseline B0 hiện tại với XLS-R + Nes2Net-X: EER(\texttt{nocodec}) = 2.71\%, EER(C04) = 41.17\%, EER(C07) = 42.32\%, \textbf{CodecGap khoảng 39.0 pp}.

\textbf{Success criterion.}
\begin{itemize}
    \item \textit{Primary}: CodecGap của B3 giảm đáng kể so với B1 và B2 (so sánh có ý nghĩa qua nhiều seed).
    \item \textit{Target}: CodecGap $< 25$ pp.
    \item \textit{Stretch}: CodecGap $< 20$ pp.
    \item \textit{Regression check}: EER trên ASV19 LA và ASV21 DF không tăng quá +1 pp so với B1 (fine-tune control). Nếu vi phạm, claim ``codec-aware training cải thiện không đánh đổi'' không đứng vững.
\end{itemize}

\textbf{Held-out codec evaluation.} Đưa vào MUST. Procedure: train với codec set $\mathcal{K}_{\text{tr}}$ (ví dụ MP3, AAC, Opus ở các bitrate đại diện); evaluate trên một dev split ASV19/ASV21 đã re-encode bằng codec set $\mathcal{K}_{\text{ho}}$ disjoint (ví dụ Vorbis, GSM, AMR-NB, Codec2, hoặc transcoding chain chưa thấy trong $\mathcal{K}_{\text{tr}}$). Báo cáo EER và CodecGap-style metric trên held-out codecs. Mục tiêu: chứng minh đề xuất cải thiện \textit{unseen-codec} generalization, không chỉ trùng C04/C07. Đây là phần làm proposal defensible nhất về novelty.

\textbf{Metric phụ trợ.}
\begin{itemize}
    \item EER tổng trên bốn dataset (ASV19 LA, ASV21 DF, ASV5 Track 1, ITW) để đo cải thiện chung và bắt regression.
    \item Grouped EER theo codec đầy đủ (C00--C11) trên ASV5 để phân tích chi tiết.
    \item ECE trên ASV5: calibration phụ trợ, không phải success criterion chính nhưng được report.
    \item Score correlation và failure overlap giữa B0 và B3 để kiểm tra đề xuất có làm thay đổi \textit{cấu trúc lỗi} (motivation cho fusion ở Phần \ref{sec:future-work}) hay chỉ scale chung.
\end{itemize}

\subsection{Rủi ro và biện pháp giảm thiểu}

\begin{itemize}
    \item \textbf{Codec mô phỏng có thể không khớp C04/C07 thực.} ASV5 không công bố cấu hình codec đầy đủ. Biện pháp: held-out codec evaluation (Phần \ref{sec:eval-protocol}) tách claim ``codec-aware generalizes'' ra khỏi ``we matched C04/C07''; ngay cả khi không trùng, held-out metric vẫn cho kết quả khoa học hợp lệ.
    \item \textbf{Consistency regularization có thể xoá cue spoof.} Biện pháp: chọn formulation score-level ($\beta$) thay vì embedding-level; sweep $\lambda$ kỹ trên dev; nếu B3 tệ hơn B2 trên ASV19 LA/ASV21 DF, hạ $\lambda$.
    \item \textbf{Fine-tune gây catastrophic forgetting trên ASV19/ITW.} Biện pháp: low learning rate cho XLS-R (partial unfreeze chỉ vài layer cuối), early stopping theo dev EER; regression check ở Phần \ref{sec:eval-protocol} bắt vấn đề này định lượng.
    \item \textbf{Compute trên Kaggle T4 hạn chế.} XLS-R 300M fine-tune cần tối ưu bộ nhớ. Biện pháp: gradient checkpointing, mixed precision, batch nhỏ + grad accumulation; nếu vẫn không đủ, giảm scope SHOULD trước, MUST sau.
    \item \textbf{Augmentation chiếm thời gian disk I/O.} Biện pháp: pre-compute codec view offline cho train set thay vì on-the-fly khi cần; tradeoff disk space vs compute.
    \item \textbf{Scope creep từ COULD.} Biện pháp: chỉ đụng đến COULD sau khi MUST có kết quả ổn định; nếu MUST không thoả success criterion, ưu tiên hiểu \textit{vì sao} trước khi mở rộng kiến trúc.
\end{itemize}

\section{Hướng phát triển} \label{sec:future-work}

Các hướng dưới đây mở ra từ Phần 3.6 và Phần \ref{sec:proposal} nhưng không thuộc core proposal; được dành cho giai đoạn sau Thực tập 2 hoặc luận văn:

\begin{itemize}
    \item \textbf{Supervised contrastive codec consistency variant.} Mở rộng formulation $\beta$ sang loss cross-utterance theo class: kéo các view \textit{bonafide} (gồm cả codec-augmented) lại gần nhau và tách khỏi cụm \textit{spoof}, và ngược lại. Variant này đòi hỏi memory bank/large batch và có rủi ro xoá cue spoof tương tự embedding-level consistency, vì vậy chỉ kích hoạt sau khi B3 score-level đã ổn định và có dev EER tham chiếu.
    \item \textbf{Score-level fusion với calibration đầy đủ (từ F4).} Sử dụng failure overlap và correlation thấp giữa các họ làm motivation cho ensemble. Yêu cầu protocol nghiêm: temperature scaling per-model trên dev split, sau đó logistic regression hoặc weighted average trên \textit{calibrated} scores, rồi eval. Báo cáo cả EER và ECE trước/sau fusion; nếu fusion giảm EER nhưng làm xấu calibration, trade-off cần được minh hoạ rõ.
    \item \textbf{Front-end / back-end ablation diện rộng (từ F5).} Mở rộng grid 2--3 front-end (XLS-R 300M, WavLM Large, multi-lingual XLS-R) $\times$ 2--3 back-end (Nes2Net-X, Nes2Net-LA local-attention variant, DuaBiMamba của XLSR-Mamba) trên cùng pipeline Phần \ref{sec:proposal}. Mục tiêu: tách phần cải thiện đến từ representation và phần đến từ thiết kế back-end.
    \item \textbf{Continual pre-training của SSL front-end trên codec-augmented audio (long-term).} Fix tận gốc F5 ở mức pre-training: tiếp tục pre-train XLS-R/WavLM trên audio codec-augmented unlabel quy mô lớn, sau đó fine-tune CM. Yêu cầu compute lớn (GPU A100 nhiều ngày), không khả thi với budget hiện tại nhưng là hướng luận văn rõ ràng.
    \item \textbf{Mở rộng benchmark.} Bổ sung MLAAD v9 (multilingual), SpeechFake (ACL 2025) và EchoFake (replay channel) sau khi verify nguồn chính thức, để đánh giá đề xuất codec-aware vượt khỏi điều kiện ASVspoof và kiểm tra tương tác giữa codec robustness và các trục robustness khác.
    \item \textbf{SASV và actual DCF.} Tính min t-DCF và a-DCF (Track 2 SASV) cùng calibration metrics đầy đủ (Brier, Cllr, actual DCF) trên các dataset có protocol phù hợp. Quan trọng cho deployment claim.
    \item \textbf{Phân tích định tính sâu hơn confident errors.} Sử dụng \codepath{hard_errors.csv} cùng spectrogram để xác định cue âm học (artifact tần cao, prosody bất thường, codec marker) mà mô hình bỏ sót; bổ sung cho phân tích định lượng của Phần \ref{sec:proposal}.
    \item \textbf{Hoàn thiện đánh giá còn thiếu.} Chạy XLS-R + AASIST trên ASV21 DF để hoàn thiện Bảng \ref{tab:eer-results}; chạy lại LFCC+LCNN với recipe mạnh hơn để có baseline hand-crafted công bằng hơn.
\end{itemize}

Đề xuất chính ở Phần \ref{sec:proposal} là khả thi với compute hiện có, có thước đo thành công định lượng (CodecGap target $<25$ pp, stretch $<20$ pp), và có held-out codec evaluation làm hàng rào defensible về novelty. Các hướng future work ở Phần \ref{sec:future-work} sẽ định hình kế hoạch dài hạn của Thực tập 2 và luận văn, được kích hoạt theo thứ tự ưu tiên sau khi proposal chính có kết quả ổn định.

% \input{chapters/Chapter5.tex}
% \input{chapters/Chapter6.tex}

\appendix
\bibliography{refs}
\bibliographystyle{unsrt}

%\bibliography{refs}
%\bibliographystyle{plain}
%-	Danh mục TL tham khảo
%-	Phụ lục (nếu có)

\end{document}
